\stepcounter{section}
\section*{Lecture 1: The animals before the taxonomy}
\begin{remark}
    An echidna is known by many different names in different languages: for example, \textit{nyingarn} (Noongar), \textit{bigibila} (Gamilaraay), \textit{tachyglossus} (scientific name), \textit{spiny anteater} (original English name). Some names are older than others; the Indigenous ones in particular have been around for tens of thousands of years. However, they all refer to the same underlying object, the animal itself, which has been around for much longer - likely tens of millions of years. Changing the name does not change the object. \\
    \\
    We shall take this as an analogy for \textbf{schemes}. The objects we call schemes are, in many cases, already familiar to us. What is new is not the objects themselves, but rather the language and perspective through which we view them. \\
    \\
    Thus, we will postpone a formal definition for this lecture. Instead, we will examine a variety of examples that can be recognised as schemes. At first, you may already recognise them through different lenses: a set, a topological space, an algebraic variety. This is much like recognising an echidna as merely an animal - such an observation is correct, but it overlooks the richer structure that distinguishes the object. Our goal through these examples is to develop an intuition for schemes before we formally define them.
\end{remark}
\begin{example}[Affine space]
    In general we can work over any field, but for now we work over $\C$, the field of complex numbers.
    \begin{enumerate}
        \item Our first example of a scheme is the set of $n$-tuples of complex numbers. This is called \textbf{affine space}, and is denoted by
        \[
        \A^n=\{(a_1,\ldots,a_n)\mid a_i\in \C\}.
        \]
        \item We can also carve out subsets of affine space to make schemes:
        \[
        X\subset\A^2, \qquad X:=\{(x,y)\mid x^2+y^2-1=0\},
        \]
        \[
        Y\subset \A^3, \qquad Y=\{(x,y,z)\mid 27x^2-y+z=0,\; 51xyz=3,\; 3z^2-5y^2=9\}.
        \]
    \end{enumerate}
\end{example}
\begin{nonexample}
    The following is \textbf{not} a scheme:
    \[
    X\subset\A^2, \qquad X=\{y^2-e^x=0\} 
    \]
    Why not? Crucially, it uses an operation that cannot be generalised to any ring: the exponential map. When we say `generalised to any ring', we mean 'uses the operations of a ring': addition, multiplication, and subtraction. This is what makes this set not an algebraic scheme - it uses something more than ring theory. 
\end{nonexample}
So far the examples given have all been schemes over $\mathbb{Z}$: the polynomials involved have all used integral coefficients, which make sense in any ring $R$ using the unique homomorphism $\mathbb{Z}\to R$. If for example, we defined the following:
\[
X\subset \A^2, \qquad X=\{(x,y)\mid ix^2-y^3=0\},
\]
this would be a scheme over $\C$, but not over $\mathbb{Z}$. That is, it makes sense in any ring that contains $\C$.
\begin{example}[Group scheme]
    We define
    \[
    \mathrm{GL}_n=\left\{\left(\begin{pmatrix}
        a_{11} &\cdots &a_{1n} \\
        \vdots & \ddots & \vdots \\
        a_{n1} &\cdots & a_{nn}
    \end{pmatrix},t\right) \Bigg|\;\; t\cdot \det(a_{ij})=1\right\}\subset \A^{n^2+1}
    \]
    This is what is called a \textbf{group scheme}. Indeed, we have maps
    \[
    \mathrm{GL}_n\times\mathrm{GL}_n\to\mathrm{GL}_n, \qquad ((a_{ij},t),(b_{ij},s))\mapsto \left(\sum_k a_{ik}b_{kj},st\right)
    \]
    \[
    \mathrm{GL}_n\to\mathrm{GL}_n, \quad (a_{ij},t)\mapsto (t\cdot \mathrm{Adj}(a_{ij}), \det(a_{ij}))
    \]
    Note that these maps are both 'algebraic': they use only ring operations. Furthermore, there exists
    \[
    I=\left(\begin{pmatrix}
        1 &  & \\
        & \ddots & \\
        & & 1
    \end{pmatrix},1\right)\in\mathrm{GL}_n.
    \]
    Together, these satisfy the group axioms, leading to the term group scheme. Other examples of group schemes include $\mathrm{SL}_n$ and $\mathrm{SO}(n)$.
\end{example}
\begin{example}[Projective space]
    We now come to a less obvious example of a scheme, called \textbf{projective space}. For now we work over $\C$. We define
    \[
    \P^n=\{\text{One dimensional (vector) subspaces of }\C^{n+1}\}
    \]
    We can view this as the set of all nonzero elements of $\C^{n+1}$, modulo the equivalence relation that
    \[
    (a_0,\ldots,a_n)\sim\lambda(b_0,\ldots,b_n), \quad \lambda\in \C\setminus\{0\}.
    \]
    That is, two points are equivalent if they differ by a nonzero scalar. We denote the equivalence class of the point $(a_0,\ldots,a_n)$ by $[a_0:\cdots:a_n]$. The space $\P^n=\{[a_0:\cdots:a_n]\}$ turns out to be a scheme.
\end{example}
\begin{example}~\label{P1=x^2-yz}
    In projective space $\P^n$, we have subschemes. For example, in $\P^2$,
    \[
    C:=\{[X:Y:Z] \mid X^2-YZ=0\}
    \]
    is a subscheme. In fact, it will turn out to be 'isomorphic' to $\P^1$ (whatever 'isomorphic' means). \\
    In $\P^5$, the following well-known space is a subscheme:
    \[
    A=\{[a:b:c:k:l:m] \mid 4k^2=2b^2+2c^2-a^2, \; 4l^2=2c^2+2a^2-b^2, \; 4m^2=2a^2+2b^2-c^2\}
    \]
    This scheme is also what is known as a K3 surface.
\end{example}
\begin{nonexample}
    Consider the subset of $\P^2$, defined by
    \[
    \{[X:Y:Z] \mid X^2-Y^2Z^2=0\}.
    \]
    A quick inspection will show that this set is not actually well-defined: $(1,1,1)$ satisfies the equation, while $(2,2,2)$ doesn't, but they both represent the same equivalence class! This shows that not every polynomial defines a well-defined subset of projective space. It turns out that only homogeneous polynomials do, as scaling will multiply the entire expression by the same factor.
\end{nonexample}
\begin{example}
    Consider the subscheme of $\P^2$, given by
    \[
    C_X:=\{[X:Y:Z] \mid X\ne 0\}
    \]
    An \textbf{inequality} is just as valid for defining a scheme as an equality is, and this is a homogeneous equation, so we are in the clear. Since $X\ne0$, we have $[X:Y:Z]=[1:y:z]$, where $y=Y/X,z=Z/X$. Since $y,z$ can take any complex value, the set $C_X$ is in bijection with $\A^2$: we have a copy of affine space lying in $\P^2$. Looking at the complement,
    \[
    \P^2\setminus C_X=\{[0:Y:Z]\},
    \]
    which is easily seen to be a copy of $\P^1$. \\
    Taking analogously $C_Y$ and $C_Z$, we see that there are three copies of $\A^2$ lying in $\P^2$. Moreover, they cover $\P^2$, since $[0:0:0]$ is not a valid element. This displays a key property of schemes: a general scheme is one which has an open cover by \textbf{affine} schemes.
\end{example}
\begin{example}[Line bundles]
    We can define the following subset of $\P^n\times \A^{n+1}$:
    \[
    \Xi = \{([a_0:\cdots:a_n],b_0,\ldots,b_n) \mid (b_0,\ldots,b_n)\text{ lies on the line spanned by }(a_0,\ldots,a_n)\}
    \]
    Equivalent formulations of this condition include:
    \begin{itemize}
        \item $\mathrm{rank}\begin{pmatrix} a_0 &\cdots &a_n \\ b_0 & \cdots & b_n \end{pmatrix}\leq 1$
        \item $a_ib_j=a_jb_i \;\;\forall i,j$
    \end{itemize}
    In this sense, $\Xi$ is an algebraic scheme. We have a natural projection map $\pi$ to $\P^n$, and we are interested in its fibers. Indeed, the points that map to $\ell\in \P^n$ are exactly those $(b_0,\ldots,b_n)$ which lie on $\ell$. In this sense, $\Xi\to \P^n$ is called a \textbf{line bundle}, as its fibers are lines. This particular one is called the \textbf{tautological} or \textbf{universal} line bundle. \\
    Another example of a line bundle is the natural projection
    \[
    \P^n\times \A^1\to\P^n.
    \]
    Here the fibers are always $\A^1$, a line. This line bundle is the \textbf{trivial} line bundle. It turns out that it is not isomorphic to the universal line bundle!
\end{example}
\begin{remark}
    The notion of 'isomorphic' has been brought up a few times now, so what do we mean when we say two schemes are isomorphic? Very generally, it means there is a bijection between them that preserves the algebraic structure: the addition, multiplication and subtraction given by the ring.
\end{remark}
\begin{example}[Algebraic isomorphism]
    Recall from Example~\ref{P1=x^2-yz} that $\P^1\cong\{X^2-YZ\}\subset\P^2$. Taking $[s:t]\in\P^1$, we have the map
    \[
    \P^1\to \{X^2-YZ\}, \qquad [s:t]\mapsto [st:s^2:t^2].
    \]
    This map is a 'morphism of schemes'. The inverse is given by
    \[
    \{X^2-YZ\}\to \P^1, \qquad [X:Y:Z]\mapsto \begin{cases}
        [Y:X] &\text{if }Y\ne0 \\
        [X:Z] &\text{if }Z\ne 0
    \end{cases} .
    \]
    Note that we never have $Y=Z=0$, since otherwise $X=0$ also. One can check that it is well-defined on the overlap $Y,Z\ne0$.
\end{example}
\begin{example}[Ring of algebraic functions]
    For this specific example, working over $\C$, let
    \[
    C=\{y^2-x^3-1=0\}\subset \A^2
    \]
    We have a natural projection to $\A^1$ given by $(x,y)\mapsto x$. We look at the ring of 'algebraic functions' on $C$ - essentially this means polynomial functions. This turns out to be
    \[
    R=\C[x,y]/(y^2-x^3-1),
    \]
    since multiples of this polynomial vanish on $C$. We define a $\C$-algebra to be a ring $R$ with a homomorphism $\C\to R$. A $\C$-algebra homomorphism $R\to \C$ by definition does nothing to $\C$, and so is determined by the images $x\mapsto a,y\mapsto b$. These must satisfy $b^2=a^3+1$ - that is, $(a,b)\in C$. As a result, points of $C$ are in bijection with $\mathrm{Hom}_\C(R,\C)$, the ring of $\C$-algebra homomorphisms $R\to \C$. \\
    Similarly considering $\A^1$, its ring of algebraic functions is $\C[x]$, and points of $\A^1$ are in bijection with $\mathrm{Hom}_\C(\C[x],\C)$, as homomorphisms are determined by the image of $x$. \\
    \\
    This bijection also translates to morphisms: given the map
    \[
    \C[x]\to \C[x,y]/(y^2-x^3-1)=R
    \]
    which sends $x\mapsto x$, we have a corresponding map
    \[
    C\cong \mathrm{Hom}(R,\C)\to \mathrm{Hom}(\C[x],\C)\cong \A^1,
    \]
    given by the projection $(a,b)\mapsto a$. That is, the morphism of schemes $C\to\A^1$ was reconstructed from the morphism of their corresponding rings, but in the opposite direction. 
\end{example}
\begin{remark}
    In time, to every ring $R$ we will associate a scheme called $\spec R$. There will be a notion of '$k$-valued points' of $\spec R$ (for $k$ a ring), and there will be a one-to-one correspondence between such $k$-valued points and $\mathrm{Hom}(R,k)$. When $k$ is an algebraically closed field, a $k$-valued point is called a \textbf{geometric point}.
\end{remark}
\begin{example}[Geometric points]
    Let $R=\Z$. We look for the $\C$-valued points of $\spec\Z$, recalling that they are in one-to-one correspondence with $\mathrm{Hom}(\Z,\C)$. There is a unique homomorphism $\Z\to\C$, and hence only one $\C$-valued point. Similarly, there is only one $\overline{\F}_p$-valued point. \\
    We now consider $\Z[i]$. There is a unique homomorphism $\Z\to\Z[i]$, and by our above discussion we should expect a corresponding scheme morphism $\spec\Z[i]\to\spec\Z$, sending $k$-valued points to $k$-valued points. To find the $\C$-valued points of $\spec\Z[i]$, we note that there are two homomorphisms $\Z[i]\to\C$, given by $i\mapsto i$ and $i\mapsto -i$. So, there are 2 distinct $\C$-valued points. For the $\overline{\F}_p$-valued points, for $p\ne2$ there are similarly $2$ homomorphisms, as there are two solutions to $x^2+1$. For $p=2$ there is only one homomorphism however. \\
    As such, our corresponding morphism $\spec\Z[i]\to\spec\Z$ must send both $\C$-valued points to the singular one in $\spec\Z$, and similarly for $\overline{\F}_p$ for $p\ne2$. For $p=2$, it send the unique $\overline{\F}_2$-valued point to the unique one in $\spec\Z$.
\end{example}


%%% Local Variables:
%%% mode: LaTeX
%%% TeX-master: "main"
%%% End:
